\uuid{Uik4}
\titre{Module/argument}
\auteur{bodin ; hanani ; mzari}
\organisation{exo7}
\licence{CC BY-NC-SA 4.0}
\datecreate{2020-01-20}
\niveau{L1}
\module{Algèbre}
\chapitre{Nombres complexes}
\sousChapitre{Racine carrée, équation du second degré}
\difficulte{1}
\theme{Équations, module, argument}
\exo7qcmid{lille-1.nombres-complexes.019}

\contenu{
\begin{question}
\qclassification{104.02,104.03,104.04}
\qtags{motcle=module/argument}
Soit $z\in \Cc$ tel que $|z-2|=1$.  Quelles sont les assertions vraies ?
\begin{answers}
    \bad{$z=3$}
    \bad{ $z=1$}

    \good{$z=2+e^{i\theta}$, $\theta \in \Rr$}

    \good{Le point du plan d'affixe $z$ appartient au cercle de rayon $1$ et de centre le point d'affixe $2$.}
\end{answers}
\begin{explanations}
$|z-2|=1$, donc $z-2=e^{i\theta}$, $\theta \in \Rr$.

\end{explanations}
\end{question}
}
